\subsection{Integral low-degree cleanup with a localized carrier}

\begin{theorem}[A-S09-31: Proved; integral mixed-cell cleanup through form three in simplex two]
\label{AS0931}
Work over $S=\Gamma_{\mathbb Z}(T)$.
Let $X\to A,B$ have torsionfree terms, $X_0=S$,
coefficient $H^0$ in every simplex and weight, and rational
coefficient comparison in every simplex and weight.
The targets are admissible. Specify a family of actual free
normalized simplex-two cells $x_i,y_i=dx_i$ of forms one and
two, weights $b_i\geq b_*$, and boundaries
$\partial x_i=dh_i$, as in AT25. Require the actual integral
presentation (AT.168), with $C$ the indicated free identity-disk
algebra on the other generators.
Choose a power $N=2^e$ and a lower bound $b_*$ such that
\[
 \begin{split}
 N&\geq\max\{1,\mu_A(0),\mu_B(0)\},\\
 b_*&\geq\max\{3,2N,\mu_A(1),\mu_B(1)\},\\
 2b_*&\geq\max_{D=A,B}\{\lambda_D(4),\mu_D(4)\}.
 \end{split}
 \tag{S09.168}
\]
For each $i$ take an actual integral graded-free resolution
of $I_1\langle2b_i\rangle$ generated in weights at least
$2b_i+1$, and form the simultaneous AT16 resolved extension
for $(\gamma_2(y_i),x_i y_i)$ and $m=2$.
Let $Y$ be its torsionfree image.

There are actual choices of integral target maps $Y\to A,B$
extending the prescribed maps on $X$, and an actual sequence
of inclusions from $Y$ whose union $Y^\sharp$ has the following
properties. Its terms are torsionfree, its rational algebra
has the coefficient comparison in every simplex and weight,
and its coefficient $H^0$ is $S$. Its first two simplices are
literally those of $X$. In simplex two it has coefficient
comparison through form three in every weight.
All extensions have actual integral maps to both targets.
After the original resolved extension, the added two-cell
repairs have cell arity at least three and annihilators of the
form $2^k$. The based degree-one saturation is integral.
The support condition $q\leq w$ is retained.

For each target there is, throughout this construction, an
actual localized carrier
$Y_2^\sharp\otimes\mathbb Z_{(2)}\to A_3\otimes\mathbb Z_{(2)}$
with zeroth face the localization of the chosen integral map
and with the reference-face bound (AT.161).
No integral carrier into $A_3$ is asserted.
The horizontal restoration of AT25, under its stated initial
two-row exactness assumption, is also available integrally and
leaves simplex two unchanged.
No vertical comparison in forms two or three in all higher
simplices, or in form four, is asserted.
\end{theorem}

\begin{proof}
We give first the two comparisons needed to pass between
integral maps and the localized carrier. Write
\[
 \mathcal C_A=\operatorname{Hom}(K(2),A)=Z_2(NA),
                    \qquad \mathcal N_A=N_3A.
\]
Both are actual cochain subcomplexes of the appropriate simplex
groups, with their coefficient actions. The normalized exact
sequences (AT.73), applied in the range (S09.168), give
$H^j(\mathcal C_A^w)=0$ for $j=0,1,3,4$ and
$H^2(\mathcal C_A^w)=S^w$ whenever
$w\geq2b_*$. The latter is the natural coefficient-period
identification. The same assertions hold after localization
at two. Exactness of normalization gives
$H^j(\mathcal N_A^w)=0$ through degree four in that range.
Horizontal exactness also gives the actual surjections
\[
       d_0:\mathcal N_A^{w,j}\longrightarrow\mathcal C_A^{w,j}
                  \quad(0\leq j\leq4,\ w\geq2b_*).
 \tag{S09.169}
\]
All these statements use the specified uniform bounds, and
hold independently for $B$.

\paragraph{Integral initial maps with the prescribed local carrier.}
The integral AT16 construction supplies initial maps on each
resolved summand. It uses the actual chosen integral resolution;
its positive-weight generators have weights at least $2b_i+1$.
After localizing this same resolution at two, AT20 and the
carrier construction of AT24 give a particular local target
map and a carrier on the simplex-two algebra. Denote that map
by $f_{\mathrm{loc}}$ and an initial integral lift by $f_0$.
Both have the same fixed map on the old algebra $X$.
We arrange their agreement at the zeroth face up to an explicitly
lifted linear homotopy, instead of identifying their choices.

Put $s=2b_i$ and $L=I_1\langle s\rangle$.
The difference of the two linear generator assignments is a
cocycle in the actual free-resolution Hom complex used in AT16.
Its degree-zero cohomology is
\[
 \operatorname{Ext}^1_{\mathrm{gr}S_{(2)}}(L_{(2)}/2,S_{(2)})
                       =S_{(2)}^s/2S_{(2)}^s.
 \tag{S09.170}
\]
Here the finite-diagonal comparison is precisely (AT.100):
all generator weights lie in the displayed period range.
The last equality follows from the multiplication-by-two
sequence and the positive source-shift comparison AT04.
This computation is on the same localized integral resolution;
no minimal local resolution is substituted.

The additional base weight $s$ lies in the period range by
(S09.168). Choose an actual integral closed
$\omega\in\mathcal C_A^{s,2}$ with period $T^{[s]}$.
The assignment $v_0(t)\mapsto\varepsilon(t)\omega$, zero
on all other generator summands, is an actual integral
cochain map on $W$, by $\varepsilon\partial=0$.
As checked in (AT.102), it represents the generator of
(S09.170). The map $S^s/2\to S_{(2)}^s/2$ is onto.
Adding an integral multiple of this assignment to $f_0$
therefore produces an integral lift $f$ such that
\[
                 f_{\mathrm{loc}}-f_{(2)}=dH+Hd_W
 \tag{S09.171}
\]
for an actual degree-minus-one $S_{(2)}$-linear homotopy $H$
on the resolved generator complex. Its degree-zero component
is zero. The terminal relation descends through $W^0$ exactly
as in the free-resolution calculation (AT.100)--(AT.102).
This is a statement about the cohomology of that Hom complex,
not an interchange of localization with an infinite product.

The modules $W^1,W^2,W^3$ are free over $S_{(2)}$.
Use (S09.169) on each homogeneous basis element to lift
$H|_{W^j}$ to
$K:W^j\to\mathcal N_{A_{(2)}}^{j-1}$ for $j=1,2,3$.
Set $K|_{W^0}=0$. If $F$ is the local carrier of AT24,
replace its values on $W$ by
\[
                              F'=F-dK-Kd_W.
 \tag{S09.172}
\]
It remains a cochain map for the same affine twist: the
change has zero twist and its differential commutes with
$d_W$. On the actual bottom module $W^0$, the term is
$-Kd_W$, so it respects that module's existing relations.
Its zeroth face is $f_{(2)}$ by (S09.171).
All positive faces of the change vanish, since $K$ takes
values in $N_3$. Consequently every reference-face value
in (AT.160), and hence the form/arity bound (AT.161),
is unchanged. On the free old algebra use the old carrier.
The free relative PD universal property supplies the resulting
actual algebra map. All integral target maps factor through
the torsionfree image. Carry out these choices independently
for each cell and each target; the simultaneous free
presentation makes them compatible on the common old algebra.
This constructs the claimed initial integral maps and local
carriers on $Y$.

\paragraph{The exact odd-prime calculation, including the bottom module.}
We next prove that $H^1(Y_2),H^2(Y_2),H^3(Y_2)$ have no
odd-primary torsion. Fix an odd prime $\ell$ and work over
$S_{(\ell)}$. More generally write $m$ for an invertible
scalar in the following coefficient-complex calculation.
For the cone (AT.75),
\[
 \begin{split}
 W^1&=Q_2u\oplus Q_1v,\qquad W^2=Q_1u\oplus Q_0v,\qquad
 W^3=Q_0u,\\
 d(a,b)&=(-\partial a+mb,\partial b),\qquad
 W^0=\ker(W^1\to W^2).
 \end{split}
\]
Invertibility of $m$ gives the exact graph identification
\[
 W^0=\{(a,\partial a/m):a\in Q_2\}\cong Q_2.
 \tag{S09.173}
\]
Thus the bottom module is free here. In $W^1$ use coordinates
$(a,c=b-\partial a/m)$ and in $W^2$ use
$(A,e=B-\partial A/m)$. The differentials become
\[
 Q_2\xrightarrow{1}Q_2,\qquad
 Q_1\xrightarrow{m}Q_1,\qquad
 Q_0\xrightarrow{m}Q_0
 \tag{S09.174}
\]
in form ranges $0\to1$, $1\to2$, and $2\to3$, respectively.
Indeed the first differential sends $a$ to $(a,0)$;
the next sends $(a,c)$ to $(mc,0)$; the last sends
$(A,e)$ to $me$. These formulas give inverse linear
changes of generators and prove the displayed decomposition.
It invokes no preservation assertion for a free PD algebra
on a nonfree module.

For the actual twist, first make the integral-over-$S_{(\ell)}$
change $u_0'=u_0-\varepsilon\eta/m$, with all other generators
unchanged. The equations $d\eta=mz$ and
$\varepsilon\partial=0$ remove every affine twist term.
The inverse adds the same old element. These are strict PD
changes by the relative free universal property; the translated
$u_0$ is odd. The linear changes (S09.173)--(S09.174) then give
an actual tensor product with free identity disks over
$S_{(\ell)}$. At $m=2$ this applies to every initial resolved
summand, including its $K(2)$ parameter, whose value in each
simplex is free.

An even-to-odd free disk has the differential identity pairs
$\gamma_{n+1}(v)\mapsto\gamma_n(v)dv$; the inverse on
each positive pair is its integral contraction.
For a free odd one-form disk its arity-$n$ complex is
\[
 S_{(\ell)}v\gamma_{n-1}(dv)
           \xrightarrow{\ n\ }S_{(\ell)}\gamma_n(dv)
 \quad\text{in forms }2n-1,2n.
 \tag{S09.175}
\]
All arities $n<\ell$ are contractible over $S_{(\ell)}$;
every remaining summand starts in degree at least
$2\ell-1\geq5$. Contracting the unit summands and the
even disks therefore leaves only the coefficient summand
and terms beginning in degree five. Tensoring with any
nonnegative old complex changes no cohomology in degrees
one through three. The contractions are $S_{(\ell)}$-linear.
Finite products and their filtered unions give the same
conclusion for all displayed free generators.
The original $X_2$ is itself the stated free identity-disk
algebra and has the same odd-prime vanishing through three.
Thus the conclusion follows for the initial resolved $Y_2$.

This comparison also identifies the torsionfree image after
localization. The localized displayed free-disk algebra is
torsionfree over $\mathbb Z_{(\ell)}$, so it already embeds
in its rationalization. Localization is exact and takes the
original torsion kernel to its torsion kernel, which is zero
here. Consequently no additional quotient is hidden in this
odd-prime calculation.

The same conclusion persists under the later repairs. For
$m=2^k$, the change $u'=u-b/m$ over $\mathbb Z_{(\ell)}$
turns a form-two repair into an even zero-form identity disk,
and a form-three repair into the odd one-form disk (S09.175).
Again the comparison through form three is unchanged.
Based saturation changes no localized zero-form in simplex
two: a permitted rational root has integral differential
there, and $H^1(Y_{2,(\ell)})=0$ supplies an integral primitive.
Their difference is a rational coefficient. Its integral
vertex value makes that coefficient integral over
$\mathbb Z_{(\ell)}$, so the root was already present.
This applies at every stage and to all PD closures.
It also covers faces of roots added in higher simplices,
since their derivative and vertex conditions are the same
allowed conditions in simplex two.

\paragraph{Every required integral repair has a compatible localized top.}
The arity calculation in AT25 is valid over $\mathbb Z$:
arity one is a sum of identity complexes; the arity-two
mixed complex has maps $(1,-1)$ and $(1,1)$; and each
diagonal has connecting map
$I_1\langle2b_i\rangle/2\hookrightarrow
S\langle2b_i\rangle/2$, which is injective in the actual
coefficient basis. Tensoring with $C$ over the PID $\mathbb Z$
introduces no degree-at-most-three term: its $H^0$ is free,
and its first positive cohomology is in degree four. Every
positive-arity complex just considered has $H^0=0$, so
the possible $\operatorname{Tor}_1$ term using $H^4(C)$
and that degree-zero cohomology is also zero. In arity two
the torsionfree image changes no linear lattice: $C$ is
free over $S$ and every weight of $W$ is free abelian.
Thus arities zero, one and two have no $H^1,H^2,H^3$.
Saturation adds no zero-form in these arities, by the same
primitive and coefficient argument. All later repairs are
assigned their actual cycle arity, so these pieces stay fixed.

Every remaining positive class in forms two or three of
simplex two is rationally zero and hence torsion.
Its localizations at every odd prime vanish by the preceding
calculation. Its finite annihilator may therefore be chosen
as $m=2^k$ for some $k\geq1$. Choose homogeneous data
$dz=0$, $db=mz$ of common arity $a\geq3$, with $z$ of
form $q=2$ or $3$. Their weights satisfy $w\geq ab_*$.
As in AT25, $b$ has all faces zero: its form $q-1<a$
excludes a monomial made only from old selected cells, and
every remaining monomial has a zero-face new-generator factor.
The equation and torsionfreeness imply the same for $z$.
Thus both are actual $K(2)$ values.
Their carrier values are in $N_3(A_{(2)})$, by (AT.161)
and the same equation. Choose local top primitives
\[
                 d\widetilde U=F(z),\qquad
                 d\widetilde V=m\widetilde U-F(b)
 \tag{S09.176}
\]
there. The fixed bounds (S09.168) suffice since
$w\geq3b_*\geq2b_*$. Let $U_{\mathrm{loc}},V_{\mathrm{loc}}$
be their zeroth faces.

Flat localization commutes with the finite kernel defining
$\mathcal C_A$. For $m=2^k$, every odd integer acts invertibly
on a module killed by $m$, giving the actual complex isomorphism
\[
 \mathcal C_A/m\mathcal C_A
        \xrightarrow{\;\cong\;}
 \mathcal C_{A_{(2)}}/m\mathcal C_{A_{(2)}}.
 \tag{S09.177}
\]
The local pair shows that $f(b)$ modulo $m$ is a boundary
in the complex on the right. Use this isomorphism to choose
$v\in\mathcal C_A^{w,q-2}$ with $dv\equiv f(b)\pmod m$.
The actual integral elements
\[
             U=(f(b)-dv)/m,\qquad V=-v
 \tag{S09.178}
\]
belong to $\mathcal C_A$, satisfy $dU=f(z)$ by
torsionfreeness, and satisfy $dV=mU-f(b)$ directly.
They define the required integral target map on the free
$K(2)$ repair $du=z,dv_{
m new}=mu-b$.

We verify that these particular integral bottom values retain
a carrier. In $\mathcal C_{A_{(2)}}$ put
$\delta U=U_{\mathrm{loc}}-U$ and
$\delta V=V_{\mathrm{loc}}-V$. Then
$d\delta U=0$ and $d\delta V=m\delta U$.
For $q=2$, $H^{q-1}(\mathcal C)=H^1(\mathcal C)=0$;
for $q=3$, it is the torsionfree coefficient group in
$H^2(\mathcal C)$. In either case $\delta U=dK$ for an
actual $K$ of degree $q-2$. The element $\delta V-mK$ is
closed, and $H^{q-2}(\mathcal C)=0$ for both choices of $q$.
For $q=3$ write it as $dH$ with $H$ in degree zero;
for $q=2$ it is already zero, and put $H=0$ in degree minus one.
Use (S09.169) to lift $K,H$ to normalized
$\widetilde K,\widetilde H\in N_3(A_{(2)})$ of the same forms.
Replace (S09.176) by
\[
 \widetilde U'=\widetilde U-d\widetilde K,\qquad
 \widetilde V'=\widetilde V-m\widetilde K-d\widetilde H.
 \tag{S09.179}
\]
These retain both differential equations, have zeroth faces
exactly $U,V$, and still have every positive face zero.
Thus they extend the localized carrier while the actual
integral maps use (S09.178). Their zero reference faces retain
(AT.161). The construction for the other target is independent
on the same integral source.

\paragraph{Saturation and the union.}
After each set of these integral extensions take its image
in its rationalization and apply A-SAT01 at cutoff
$2b_*+1$. Every new weight is at least $3b_*$, and all
positive-arity zero-form roots in simplex two have zero
faces. The argument of AT24 applies to the localized carrier:
a rational root has integral local derivative, its image has
an integral local primitive, and its zero vertex value removes
the rational coefficient difference. Its positive reference
faces are zero by (AT.161) applied to its derivative.
Thus the carrier extends through the integral saturation too.
The actual integral target maps extend by A-SAT01 itself.
The homogeneous arity decomposition and its low-arity fixedness
are exactly the integral versions just verified above.

Every closed one-form in simplex two that can be new has arity
at least three and weight at least $3b_*$. Hence it lies in
the saturation range. The odd-prime calculation shows at each
stage that all form-two and form-three classes still admit
$2$-power annihilators. Iterate the operation on every such
homogeneous class and saturate after each stage. Each element
uses finitely many generators. In the countable union, every
closed two-form or three-form lies in a stage and its class
is killed at the next repair stage; closed one-forms are
handled by the next saturation stage. This proves the claimed
coefficient comparison in simplex two in every weight.

Rational comparison in every simplex and weight is retained
by the actual change $u'=u-b/m$ and the rational monomial
disk contractions; saturation does not change rationalization.
Every new integral closed zero-form is a rational coefficient
whose image in a target is integral in the stipulated high
weight range, so coefficient $H^0$ is retained. All terms stay
torsionfree by construction. Every $K(2)$ attachment fixes
simplices zero and one. To check saturation there as well,
$X_1$ is a cochain retract of the free identity-disk algebra
$X_2$, by $d_0s_0=1$. That algebra has $H^1=0$, so
$H^1(X_1)=0$. Together with its coefficient $H^0$ and the
integral vertex condition this makes every allowed rational
zero-form already integral in $X_1$. Thus the first skeleton
is fixed throughout. The form and weight bounds for the
original and new generators prove the support assertion.

Finally the optional horizontal restoration is the actual
C21 construction in dimensions at least three from AT25.
Fixing simplices zero and one while enlarging simplex two
retains the old row-zero homology and makes row-one homology
a quotient of the old one. Its given eventual vanishing
therefore persists. Higher-dimensional C21 cells restore the
remaining rows with actual target maps and leave simplex two
unchanged. This proves the precisely stated optional assertion
and completes the integral construction.
\end{proof}
